\section{Related work}\label{sec:related}

\paragraph{Bias and its evaluation in scene graph generation.}\label{sec:related-applied}
Scene graph generation was posed as structured prediction over a whole image, solved by
message passing between object and relation nodes~\cite{xu2017scene}; MOTIFS then showed a
frequency predictor of the predicate given the object pair to be surprisingly
competitive~\cite{zellers2018neural}. Mean recall became the headline metric, and debiasing methods now span the causal TDE
intervention~\cite{tang2020unbiased}, loss re-weighting and group-wise
learning~\cite{desai2021learning,dong2022stacked}, predicate-rebalancing message
passing~\cite{li2021bgnn}, relabelling of the training data~\cite{zhang2022finegrained}
and, outside scene graphs, post-hoc logit adjustment~\cite{menon2021long}; the line
continues~\cite{kuang2024debiasing,luo2025rebalancing}, and the metrics themselves have
been shown to reward pathologies~\cite{li2022rethinking}, a diagnosis that recurs in
open-vocabulary settings~\cite{li2024pixels}. Benchmark audits outside scene graphs find the
same pattern: VQA models collapse when answer priors change between train and
test~\cite{agrawal2018dont}, GQA-OOD stratifies accuracy by answer rarity within question
groups~\cite{kervadec2021roses}, and multimodal models exploit non-visual
shortcuts~\cite{geirhos2020shortcut,zhou2025shortcut}. Gains can also be credited to a
mechanism that did not produce them: grounding-based VQA methods keep their
out-of-distribution improvement when the grounding signal is
irrelevant~\cite{shrestha2020negative}, and methods can exploit how an
out-of-distribution split was built~\cite{teney2020value}. Long-tailed recognition faces the
same question for re-weighting and logit adjustment~\cite{cao2019learning,menon2021long}.
All of these characterise a \emph{single} model. We attribute the gain \emph{between} two.

\paragraph{Proper-score decompositions.} Conditional on a grouping, a proper score splits
into reliability and resolution terms~\cite{murphy1973vector,degroot1983comparison}, for
general proper scores~\cite{brocker2009reliability}. Our case-level term is not that
resolution term: it descends from the \emph{covariance} decomposition of the Brier score
due to Yates~\cite{yates1982external}, which a monotone recalibration moves while the
classical resolution is invariant to it. The two coincide only when both models are
calibrated within every group. Variance estimators exist for single-model
decompositions~\cite{siegert2014variance}, but not for a contrast between two models
scored on the same data. The grouping loss~\cite{perezlebel2023beyond} is the
single-model relative of our case-level part: what calibration leaves behind because
inputs sharing a confidence level differ; ours contrasts two models within groups
declared in advance.

\paragraph{Comparing two predictors.} Inference on a paired score difference is the
Diebold--Mariano test~\cite{diebold1995comparing} and its conditional
form~\cite{giacomini2006tests}; at a single group our total-gain interval is a clustered
Diebold--Mariano test (\cref{cor:dm}), so the decomposition adds a split to what those tests
already report. DeLong \etal~\cite{delong1988comparing} compare correlated ROC areas with a
paired U-statistic, the same device we use for a group-conditional score gain. The
case-level term is an order-two U-statistic~\cite{hoeffding1948class}, with a H\'ajek
projection supplying its influence function~\cite{serfling1980approximation}.
